% Part: lambda-calculus
% Chapter: introduction
% Section: lambda-definability

\documentclass[../../../include/open-logic-section]{subfiles}

\begin{document}

\olfileid{lam}{int}{rep}
\olsection{\usetoken{S}{lambda definable} అంకగణిత ప్రమేయాలు}

లాంబ్డా కలనశాస్త్రం గణనా నమూనాగా ఎలా పనిచేయగలదు? మొదట చూసినప్పుడు,
ఈ ప్రశ్నకు అర్థం ఎలా ఇవ్వాలో కూడా స్పష్టంగా లేదు. సహజ సంఖ్యలపై
గణనీయత గురించి మాట్లాడాలంటే, ఆ సంఖ్యలకు తగిన ప్రాతినిధ్యాన్ని కనుగొనాలి.
ఆశ్చర్యకరంగా బాగా పనిచేసే ఒక ప్రాతినిధ్యం ఇదిగో.

\begin{defn}
ప్రతి సహజ సంఖ్య~$n$కు, \emph{చర్చ్ సంఖ్యాంకం} $\num{n}$ను లాంబ్డా
పదం $\lambd[x][\lambd[y][(x(x(x(\dots
x(y)))))]]$గా నిర్వచించండి; ఇందులో మొత్తం $n$ $x$లు ఉంటాయి.
\end{defn}

$\num{n}$ పదాలు ``పునరావర్తకాలు'': $f$ను ఇన్‌పుట్‌గా ఇస్తే,
$y$ను $f^n(y)$కు పంపే ప్రమేయాన్ని $\num{n}$ తిరిగి ఇస్తుంది. ప్రతి
సంఖ్యాంకం నియత రూపంలో ఉందని గమనించండి. సహజ సంఖ్యలపై ఒక ప్రమేయాన్ని
లాంబ్డా పదం ``గణించడం'' అంటే ఏమిటో ఇప్పుడు చెప్పవచ్చు.

\begin{defn}
$f(x_0, \dots, x_{k-1})$ అనేది $\Nat^k$ నుండి $\Nat$కు వెళ్లే
$k$-స్థానిక పాక్షిక ప్రమేయం అనుకుందాం. ప్రతి సహజ సంఖ్యల అనుక్రమం
$n_0$, \dots,~$n_{k-1}$కు, $f(n_0, \dots, n_{k-1})$ నిర్వచితమైతే
\[
F\, \num{n_0}\, \num{n_1} \dots \num{n_{k-1}} \red \num{f(n_0, n_1, \dots,
  n_{k-1})}
\]
అవుతూ, లేకపోతే $F\, \num{n_0}\, \num{n_1} \dots \num{n_{k-1}}$కు
నియత రూపం లేకపోతేనే, $\lambd$-పదం~$F$ ప్రమేయం~$f$ను
\tetoken{లాంబ్డాతో నిర్వచిస్తుంది}{!!{lambda define}s} అని అంటాం.
\sourcecorrection{OLTELAMCOMP-001}{ఆంగ్ల మూలం $k$ ఆర్గ్యుమెంట్లు గల $f(x_0,\dots,x_{k-1})$ను $n$-స్థానికమనీ, $\Nat$ నుండి $\Nat$కు వెళ్లేదనీ పేర్కొంది. వెంటనే వచ్చే $k$ ఇన్‌పుట్‌ల నిర్వచనానికి అనుగుణంగా దాన్ని $\Nat^k$ నుండి $\Nat$కు వెళ్లే $k$-స్థానిక పాక్షిక ప్రమేయంగా పునరుద్ధరించాం.}
\sourcecorrection{OLTELAMCOMP-002}{అనిర్వచిత సందర్భంలో ఆంగ్ల మూలం $F$ తరువాత ప్రయోగ ఖాళీకి బదులు కామాను ఉంచింది. నిర్వచిత సందర్భంలోని అదే $k$-ఆర్గ్యుమెంట్ల ప్రయోగ పదానికి నియత రూపం లేదనే ఉద్దేశిత షరతును పునరుద్ధరించాం.}
\end{defn}

\begin{thm}
\ollabel{thm:lambda-def}
ప్రమేయం~$f$ పాక్షిక గణనీయ ప్రమేయమైతేనే, అప్పుడే మాత్రమే అది ఒక
లాంబ్డా పదంచే \tetoken{లాంబ్డాతో నిర్వచించబడిన}{!!{lambda defined}} అవుతుంది.
\end{thm}

\begin{explain}
ఈ సిద్ధాంతం కొంత ఆశ్చర్యకరమైనది. గణనా నమూనాగా లాంబ్డా కలనశాస్త్రం
చాలా సరళమైన కలనశాస్త్రం; ఇందులో లాంబ్డా అమూర్తీకరణ, ప్రయోగం మాత్రమే
క్రియలు!{} అయినప్పటికీ, ఇంత స్వల్పమైన వనరులతో ఏ గణనా ప్రక్రియనైనా
అమలు చేయవచ్చు.
\end{explain}

\end{document}
